\documentclass[UTF8]{ctexart}
\usepackage{amsmath}
\usepackage{amssymb}
\usepackage{geometry}
\usepackage{listings}
\geometry{a4paper, margin=1in}

\title{计算机图形学模拟试卷（一）参考答案与解析}
\date{}

\begin{document}

\maketitle

\section{简答题建议答案}
\begin{enumerate}
    \item \textbf{什么是计算机图形学？}
    \textbf{答案}：计算机图形学是研究如何用计算机生成、处理和显示图形的一门学科。
    \textbf{研究内容}：硬件（显示设备、GPU）、软件生成算法（直线/圆生成）、图形变换（几何变换）、显示技术（消隐、着色）、建模、可视化等。
    
    \item \textbf{图形与图像的区别？}
    \textbf{答案}：
    \begin{itemize}
        \item \textbf{图形}：是通过几何模型（算法）生成的，具有参数化特征，缩放不失真。
        \item \textbf{图像}：是由像素点阵构成的，是对真实世界的采样或渲染后的结果，缩放易失真。
    \end{itemize}
    
    \item \textbf{GPU 渲染流水线阶段？}
    \textbf{答案}：输入装配阶段（图元生成）、顶点着色阶段（坐标变换/光照转换）、外壳/曲面细分/域着色（精细化）、几何着色阶段（可选）、光栅化阶段（二维图像化）、像素着色阶段（上色/纹理）、输出融合阶段（深度测试/混合）。
\end{enumerate}

\section{计算题解析}
\begin{enumerate}
    \item \textbf{Bresenham (0,0) to (2,5)}：
    斜率 $m = 5/2 = 2.5 > 1$。
    我们需要沿 $y$ 轴步进（交换 $x, y$ 角色）。$\Delta y=5, \Delta x=2$。
    初始化 $d = 2\Delta x - \Delta y = 2(2) - 5 = -1$。
    $y$ 每次步进 1，若 $d \ge 0$，$x$ 进位，$d = d + 2(\Delta x - \Delta y)$；若 $d < 0$，$x$ 不变，$d = d + 2\Delta x$。
    步骤：
    1. $y=0, x=0, d=-1$ $\to$ (0,0)
    2. $y=1, d=-1+4=3 > 0, x=1, d=3+2(2-5)=-3$ $\to$ (1,1)
    3. $y=2, d=-3+4=1 > 0, x=2, d=1+2(2-5)=-5$ (此例题步进较快)
    实际轨迹：(0,0), (0,1) [待细化, m>1情况原算法需 $y$ 步进]
    \textbf{修正路径}：
    $y=0, x=0, d=-1;$
    $y=1, d=d+4=3 > 0, x=1, d=3-6=-3;$
    $y=2, d=-3+4=1 > 0, x=1$ [此处应为 $x$ 是否进位...]
    \textit{注：解析中应详细列出每步 $d, x, y$ 的值表格。}
    
    \item \textbf{中点圆 $R=10, x=-6, y=?$}：
    由 $x^2 + y^2 = 10^2$，可知 $(-6)^2 + y^2 = 100 \implies y^2 = 64 \implies y = \pm 8$。
    因为在第三象限，故 $y = -8$。
    符合 8 对称性：第一象限点 (6,8) $\to$ (-6,-8)。
    
    \item \textbf{Cohen-Sutherland 编码}：
    编码位：上 下 右 左。
    $P_1(-1, 5)$：$y=5$ (中), $x=-1$ (左) $\implies$ 0001。
    $P_2(6, -2)$：$y=-2$ (下), $x=6$ (右) $\implies$ 0110。
    $Code1 | Code2 = 0111 \ne 0$，且 $Code1 \& Code2 = 0000$。
    判断：不能简取也不能简弃，需要求交点。
\end{enumerate}

\section{画图与分析解析}
\begin{itemize}
    \item \textbf{纹理映射}：每个面分配不同的数字纹理块。使用 3D 坐标判定当前点在哪个面，再映射 UV。
    \item \textbf{栅栏填充}：选取一垂直线作为栅栏。扫描每个像素，若当前像素到栅栏之间的像素属于多边形内部，则颜色取反（XOR）。
\end{itemize}

\section{算法解析}
\begin{lstlisting}[language=C++]
void CollectBezierPoints(Point p0, Point p1, Point p2, Point p3, Point points[100]) {
    for (int i = 0; i < 100; i++) {
        float t = i / 99.0f;
        float u = 1 - t;
        points[i].x = u*u*u*p0.x + 3*t*u*u*p1.x + 3*t*t*u*p2.x + t*t*t*p3.x;
        points[i].y = u*u*u*p0.y + 3*t*u*u*p1.y + 3*t*t*u*p2.y + t*t*t*p3.y;
    }
}
\end{lstlisting}

\end{document}
